\section{League 训练机制：exploiter、不对称博弈与匹配策略}

前面的失败谱系说明单一对手无法持续暴露策略盲区；本节进一步回答如何组织一整个训练 population。League 的重点不在并行模型数量，而在不同角色承担互补目标，并由 match-maker 选择既有挑战性又能产生学习信号的对局。

在 44--53 分钟段，讲者详细讲了 AlphaStar league 如何避免策略塌缩。核心机制不是简单``多放几个模型一起打''，而是通过不对称训练目标和精细 match-making 维持压力。

\begin{figure}[H]
\centering
\includegraphics[width=0.92\textwidth]{frame-04.jpg}
\caption{Self-play 与引入 exploiter 后的策略演化差异}
\label{fig:selfplay-vs-exploiter}
\end{figure}

\begin{knowledgebox}{读图：Self-play 与 exploiter 的压力不同}
纯 self-play 的双方共同演化，可能围绕当前习惯形成狭窄均衡；exploiter 不追求长期通用性，而专门寻找 main policy 的局部漏洞。图中的分支扩展表示策略覆盖增加，不代表每个 exploiter 都应直接部署。其价值是生成训练信号与回归场景。
\end{knowledgebox}

\begin{figure}[H]
\centering
\includegraphics[width=0.95\textwidth]{diagrams/league-training.png}
\caption{讲义整理图：League 中的 main policy、exploiters、历史策略、specialists 与 match-maker}
\label{fig:league-training}
\end{figure}

\begin{knowledgebox}{概念解释：League 五个角色如何互补}
main policy 优化部署目标；exploiters 寻找当前弱点；historical policies 防止遗忘旧漏洞；specialists 压测少见策略；match-maker 根据胜率、不确定性和新颖度安排对局。archive 保存 checkpoint、trace 与评分，使训练压力可以复现，而不是由临时 prompt 随机决定。
\end{knowledgebox}
图 \ref{fig:selfplay-vs-exploiter} 对比了两条演化轨迹\footnote{来源：\href{https://www.youtube.com/watch?v=ntjOxjZMaac}{Lecture 03 视频帧}，时间戳 00:44:50。}：仅自博弈容易走向狭窄专精，而加入 exploiter 后会被迫寻找更鲁棒策略。

\begin{importantbox}{exploiter 的关键设计}
讲者描述了一个重要不对称：主策略（main agent）不总能看到 exploiter 的全部训练信息，而 exploiter 专注于击败主策略。这个不对称会持续暴露主策略盲点，抑制过早收敛。
\end{importantbox}

\begin{knowledgebox}{match-making 的两个候选策略}
\begin{enumerate}
    \item 对 hardest opponent 训练：直接面对最强压制者，但可能学习信号过稀。
    \item 对 50/50 opponent 训练：接近最大熵区间，梯度信号更稳定。
\end{enumerate}
字幕中讲者指出 50/50 常带来更平衡的方差与学习信号（00:52:35--00:52:54）。
\end{knowledgebox}

我们可以用信息论语言表达其直觉：若胜率 $p$ 接近 0 或 1，策略更新常退化；当 $p\approx0.5$ 时，不确定性最高，改进空间最大。二元熵
\[
H(p)=-p\log p-(1-p)\log(1-p)
\]
在 $p=0.5$ 取最大值，这与讲者的经验观察一致。

\begin{warningbox}{只追 hardest 不一定最优}
如果对手强到长期碾压，模型拿不到有效学习信号；如果对手弱到长期碾压对方，也学不到新东西。match-making 本质是在``挑战性''与``可学习性''之间找平衡。
\end{warningbox}

\subsection{本章小结}
League 的本质是组织学习压力，而不仅是增加算力。exploiter 机制、对手采样和训练不对称性共同决定了策略最终形态，这些思想可直接迁移到 LLM Agent 的 red teaming 与在线优化流程。

